\documentclass{subfiles}
\begin{document}
    Defined starting from gamma encoding,
        delta encoding allows to define shorter codewords\footnotemark.
        By using the same notation as before, if \(x\) is an integer 
        its gamma encoding consists of two parts: the gamma encoding of 
        \(\abs{B(x)}\) and \(\tilde{B}(x)\). That is
        \[
            \delta(x) = \gamma(\abs{B(x)}) \tilde{B}(x).
        \]

    \footnotetext{It should be pointed out that this is true for values of \(n \ge 32\).}

    \begin{example*}
        Consider again 29. Its binary representaion is \(B(29) = 11101\).
            The gamma encoding of \(\abs{B(29)}\) is 
            \[\begin{aligned}
                \gamma(29) & = U(\abs{B(\abs{B(29)})}) \tilde{B}(\abs{B(29)}) \\ 
                    & = U(\abs{B(5)}) \tilde{B}(5) \\ 
                    & = 001 01.
            \end{aligned}\]
            
            We can conclude that the delta encoding of \(29\) is 
            \[
                \delta(29) = \underbrace{
                        \overbrace{001}^{U(\abs{B(5)})}
                        \overbrace{01}^{\tilde{B}(5)}
                    }_{\gamma({\abs{B(29)}})}
                    \underbrace{1101}_{\tilde{B}(29)}
            \]
    \end{example*}

    Decoding is simple: count the number of zeros up to the first 1,
        say these are \(k\), fetch the following \(k + 1\) and interpret the sequence 
        as the integer \(x\); fetch the following \(x - 1\) bits, 
        prepend a 1 to the sequence and interpret it as the integer \(n\).

    It can be proved that delta encoding is optimal when the distribution is 
    \[
        p(x) = \frac{1}{x \log^{2} x}.
    \]

    \begin{remark*}
        What can we say about the sensitivity of gamma and delta encoding?
            That is, how well do the two encoding do when errors occur in transmission?
            It's easily seen that, even if a single gets flipped, the whole sequence 
            gets decoded incorrectly.
        
        The main issue is that neither of the two encoding produces self-syncronizing codewords.
            In the next section we describe a code that instead does produce self-syncronizing
            codewords

    \end{remark*}
\end{document}
